\section{Discussion}

\subsection{Principal Findings}

This study formulates tissue-specific presentation preference by a \plainMHCI{} molecule as a problem in which one model learns several related biological tasks from \plainMatchedPairs{} (paired multi-task learning). Rather than asking whether a peptide can bind to or be presented by a \plainMHC{} molecule in general, the benchmark asks which of two peptides sharing \plainMHC{} restriction and parent-protein evidence (matched peptides) is more strongly associated with a particular tissue--antigen-presenting-molecule task (tissue--MHC task). This formulation controls two major sources of trivial discrimination and focuses evaluation on relative evidence conditional on tissue (tissue-conditioned evidence).

Four findings carry the main biological message. First, peptide sequence contains reproducible tissue-associated information after the two peptides are required to share their \plainMHC{} restriction and source protein (matched on MHC and source protein). Second, the finding appears in both human and mouse, although the species use different models and should not be compared as a controlled species experiment. Third, performance falls substantially when a test peptide is absent from every training tissue--antigen-presenting-molecule combination (training task), showing that repeated peptide identities make the ordinary benchmark easier. Fourth, removing tissue identity causes a further large decrease, whereas predictors of general \plainMHC{} binding and presentation recover only part of the signal. Tissue context therefore matters to the ranking problem, but the current sequence-only model cannot identify whether the signal comes from tissue expression, antigen processing, cell composition, or study structure.

The results from familiar tissue--antigen-presenting-molecule combinations with peptide reuse across combinations allowed (standard evaluation) and from \plainPeptideOOF{} answer complementary questions. The ordinary benchmark measures performance on new peptide pairs sharing restriction and parent protein (matched pairs) within familiar combinations (tasks). The harder benchmark measures exact peptide identities absent from all training combinations within those same biological tasks. Neither evaluates a new tissue, a new allele, or an independent biological cohort.

\subsection{Effects of Model Structure under Peptide Separation}

Differences between model structures (architecture effects) are more limited under \plainPeptideOOF{} than under the ordinary benchmark. In human, TissuePMHC remains above simple sequence references, \plainSharedHeads{}, and the \plainPlainDual{}, but it is statistically comparable to the simpler \plainAuxDual{}, which also retains tissue and \plainHLA{} grouping information during training. Additional tissue and allele guidance during training (auxiliary supervision) is therefore the most reproducible retained component benefit; this comparison does not isolate superiority of the peptide encoder that examines several motif lengths (multi-scale peptide encoder) itself.

The mouse interpretation is more conservative. The \plainMMoE{}, \plainSharedHeads{}, and the \plainBLOSUM{} random forest perform similarly under peptide separation. Averaging predictions across several \plainSeeds{} contributes more strongly than the selected model structure. The mouse result therefore supports that the biological prediction problem can be learned (task learnability), rather than a universal advantage of one model structure.

\subsection{Relationship to General Peptide--Antigen-presenting-molecule Prediction}

TissuePMHC is not a replacement for a predictor of general peptide--antigen-presenting-molecule compatibility (peptide--MHC predictor). General tools ask whether a peptide is compatible with a \plainMHC{} molecule or likely to be presented overall. TissuePMHC asks whether it is preferred in one tissue relative to a comparison peptide sharing presenting restriction and parent protein (matched peptide). MHCflurry and NetMHCpan achieve AUROCs of about 0.65--0.69 on this task, confirming that general presentation propensity is relevant. Their lower performance, together with the \plainMHCOnly{} models, shows that general compatibility does not fully explain the ranking conditional on tissue.

This distinction also limits causal interpretation. Tissue-associated sequence patterns may reflect source-protein expression, antigen processing, cell composition, experimental sampling, or residual study structure. Because the model receives only peptide sequence and an identifier for the tissue--antigen-presenting-molecule combination (task label), it cannot identify which process produced an association. Adding a general presentation score changes AUROC by at most 0.0041, suggesting that little additional generic signal remains after the tissue-conditioned score is known.

\subsection{Potential Applications}

The benchmark may support several research applications. Tissue-conditioned scores could help prioritize candidate antigens whose presentation evidence is concentrated in a target tissue, rank potential normal-tissue off-targets during early therapeutic screening, and select peptides for follow-up immunopeptidomics experiments. The same framework could also contribute to candidate triage for T-cell receptor therapies, vaccines, and other immunotherapies when used alongside binding affinity, immunogenicity, expression, safety, and population-coverage evidence.

These are potential uses rather than validated clinical applications. The balanced evaluation tasks do not reproduce clinical prevalence, and retrospective database associations cannot establish safety or therapeutic benefit. Deployment would require prospective validation, calibrated decision thresholds, uncertainty estimates, and cohorts representative of the intended use. In particular, a low predicted score must not be treated as evidence that a peptide is absent from a normal tissue.

\subsection{Limitations}

The benchmark has several important limitations. The negative member of a pair is a constructed negative based on missing database evidence (pseudo-negative), not an experimentally confirmed non-presented peptide. Incomplete tissue sampling and limited detection depth can therefore convert unobserved positives into apparent negatives. This issue is especially relevant for peptides reported across multiple tissues and for tissues with heterogeneous cellular composition.

The source data are observational and inherit biases from studies in the \plainIEDB{}. Tissue coverage, \plainHLA{} frequency, sample preparation, mass-spectrometry platform, reporting practice, and detection depth are uneven. Matching on the \plainMHC{} and parent protein reduces specific confounders but does not remove study, laboratory, donor, or experimental-batch structure. The deliberately balanced 1:1 test tasks are useful for controlled ranking comparisons but do not represent real-world prevalence; consequently, benchmark AUPRC and score thresholds should not be transferred directly to deployment settings.

Source-tissue text is retained exactly as recorded in the \plainIEDB{}. Synonyms, anatomical subregions, and differently formatted labels are not mapped to a common controlled biological vocabulary (ontology-normalized), so semantically related tissues may be split into separate tasks. Resolving those labels would change task definitions and require rebuilding the benchmark and retraining all models; the present revision therefore documents the issue rather than mixing incompatible results.

The present work covers unmodified canonical 9-mer peptides presented by a \plainMHCI{} molecule. It does not address variable-length \plainMHCI{} ligands, post-translationally modified peptides, spliced peptides, or presentation by a class II antigen-presenting molecule (MHC-II). It also excludes explicit tissue transcriptomics, proteomics, source-protein abundance, flanking sequence, cell-type composition, and antigen-processing features. Tissue information is represented through the task definition and learned rules for sharing model parameters (sharing structure), rather than through direct biological measurements.

Generalization is restricted in several ways. Under evaluation on familiar tasks with cross-task information sharing (standard evaluation), peptides and parent proteins can recur across tasks and across fitting and evaluation data. Even \plainPeptideOOF{} permits recurrence of parent proteins, studies, platforms, and similar sequence motifs. It is stricter with respect to peptide identity but does not enforce new parent proteins (\plainProteinDisjoint{}), new studies (study-disjoint), future data (temporal split), an external cohort, or fully independent biological validation. In addition, separate output layers for each biological task (task-specific output heads) mean that all protocols evaluate previously represented tissue--antigen-presenting-molecule tasks (tissue--MHC tasks); they do not test an unseen tissue, an unseen antigen-presenting-molecule allele (MHC allele), or an unseen tissue--antigen-presenting-molecule combination (tissue--MHC combination).

The human \plainFixedTest{} was used repeatedly during a long model-development process. Although final comparisons follow a protocol that was \plainFrozen{}, project-level selection may have adapted modeling decisions to this benchmark and may produce optimism not captured by within-run confidence intervals. The mouse benchmark uses a separate internal \plainFixedTest{} that was \plainFrozen{}, but 81.47\% of its unique test peptides and 88.88\% of its unique parent identifiers from the UniProt protein database appear in training. It must not be described as preventing all entity overlap (entity-disjoint) or as external validation.

MHCflurry and NetMHCpan are \plainGeneralPMHC{} whose external scores were \plainFrozen{}, not competitors guaranteed to be free of overlap with their own training data. Their pretraining may include observations from the \plainIEDB{} or related immunopeptidomic studies, and \plainPeptideOOF{} prevents overlap only with the fitting data used by the present tissue-conditioned models. It cannot guarantee that held-out peptides were absent from external pretraining collections. In addition, the near-tie between MHCflurry and the NetMHCpan \plainEL{} in human is descriptive; without a prespecified direct paired model-comparison test, it does not establish that either external predictor is superior.

Finally, the selected main model structures differ: TissuePMHC is used for human and \plainMMoE{} for mouse. Because benchmark scale, \plainMHC{} system, and data composition also differ, the cross-species results cannot identify species effects. Within each species, comparisons on identical \plainPeptideOOF{} folds further show that superiority of one model structure is not universal: the strongest human \plainAuxDual{} is comparable to the full human model, and all three mouse models whose predictions are averaged across \plainSeeds{} are comparable.

No control in which tissue labels or pathway assignments are randomly reassigned (shuffled-tissue-label or shuffled-branch control) is reported. Such a control requires refitting models after breaking the tissue association and cannot be reconstructed from archived predictions alone. The \plainMHCOnly{} models answer the narrower completed question of how well the same model families perform when tissue information is absent.

\subsection{Future Work}

Several extensions follow directly from these limitations. First, splits that prevent recurrence of parent proteins (\plainProteinDisjoint{}), prevent recurrence of studies (study-disjoint), separate future from past data (temporal), or use external cohorts should progressively test whether performance survives the removal of additional overlap and shortcuts based on data origin (provenance shortcuts). Refitting after randomly reassigning tissue labels or model pathways (shuffled-tissue and shuffled-branch refits) should test whether the observed tissue-associated gain exceeds the gain available solely from the way tissue--antigen-presenting-molecule combinations are defined and shared (task structure). Auditable checks against external predictor training collections would also strengthen the interpretation of \plainGeneralPMHC{} whose scores were \plainFrozen{}.

Second, the input representation can be expanded to variable-length peptides and ligands presented by class II antigen-presenting molecules (MHC-II ligands). Incorporating tissue transcriptomics, proteomics, source-protein abundance, flanking sequences, proteasomal cleavage features, and cell-type composition could help separate peptide-intrinsic patterns from tissue context. Such additions should be evaluated incrementally under splits designed to prevent inappropriate overlap (leakage-aware splits), so that apparent biological gains are not driven by shared studies or entities.

Third, methods that treat recorded examples as positive but do not assume every unrecorded example is negative (positive--unlabeled learning), together with explicit models of whether an observation is likely to be recorded, may better reflect the fact that database absence is not biological absence. Calibrated uncertainty estimates could identify tissue--antigen-presenting-molecule combinations (tasks) and peptides for which the evidence is insufficient. Creating representations of biological combinations from tissue and \plainMHC{} features (conditional task representations), rather than relying only on a separate output layer for each combination (task-specific head), could enable principled evaluation on unseen tissues, alleles, or tissue--antigen-presenting-molecule combinations (tissue--MHC combinations).

Finally, controlled cross-species pretraining and adaptation could test whether learned features transfer beyond the shared definition of the biological prediction problem (shared task construction). Such experiments require harmonized data and baseline models evaluated under directly comparable conditions (matched baselines).
